\section{工程实现与训练设置}

\subsection{配置驱动的运行流程}

项目以 \texttt{src/main.py} 作为统一入口。一次正式运行依次完成配置继承、随机种子设置、数据集构建、模型初始化、训练与验证、检查点保存以及多协议测试。公共参数写在 \texttt{configs/base.yaml}，模型配置只覆盖 \texttt{run\_name}、训练模式和结构差异，减少多份 YAML 之间的隐性偏差。

每次实验写入独立目录 \texttt{outputs/<run\_name>/}，其中包含 \texttt{checkpoints/}、\texttt{logs/}、\texttt{tables/} 和 \texttt{figures/}。正式流程被拆成三个阶段：首先训练六种配置在五个随机种子下的 30 个模型并完成基础评估；随后加载各自的 \texttt{best.pt}，执行耗时较长的位置网格和固定角度扫描；最后由 \texttt{analyze\_results.py} 生成随机种子 42 的细节图，再由 \texttt{analyze\_multiseed\_results.py} 汇总五个随机种子的表格与统计图。图\ref{fig:experiment-pipeline}概括了这一过程。三个阶段可以独立重跑，因此修改报告图片不需要重新训练模型。

\begin{figure}[H]
  \centering
  \begin{tikzpicture}[
    node distance=6mm and 5mm,
    stage/.style={rounded corners=3pt, draw, line width=0.75pt, align=center,
      minimum height=13mm, text width=28mm, inner sep=3pt, font=\small},
    arrow/.style={-{Stealth[length=2.5mm]}, line width=1pt, draw=PVBlue},
    note/.style={font=\scriptsize, text=PVMuted, align=center}
  ]
    \node[stage, draw=PVBlue, fill=PVBlueLight] (config) {YAML 配置\\数据集与模型};
    \node[stage, draw=PVBlue, fill=PVBlueLight, right=of config] (train) {训练循环\\loss / accuracy};
    \node[stage, draw=PVGold, fill=PVGoldLight, right=of train] (valid) {中心验证\\早停与选优};
    \node[stage, draw=PVGreen, fill=PVGreenLight, right=of valid] (best) {最佳状态\\\texttt{best.pt}};

    \node[stage, draw=PVTeal, fill=PVTealLight, below=12mm of best] (protocol) {测试协议\\平移 / 网格 / 角度};
    \node[stage, draw=PVTeal, fill=PVTealLight, left=of protocol] (metrics) {指标与记录\\JSON / CSV / 预测};
    \node[stage, draw=PVGold, fill=PVGoldLight, left=of metrics] (analysis) {结果汇总\\表格与跨模型图};
    \node[stage, draw=PVGreen, fill=PVGreenLight, left=of analysis] (report) {报告输出\\图片与 LaTeX 表};

    \draw[arrow] (config) -- (train);
    \draw[arrow] (train) -- (valid);
    \draw[arrow] (valid) -- (best);
    \draw[arrow] (best) -- node[right, note]{加载对应\\检查点} (protocol);
    \draw[arrow] (protocol) -- (metrics);
    \draw[arrow] (metrics) -- (analysis);
    \draw[arrow] (analysis) -- (report);

    \begin{scope}[on background layer]
      \node[rounded corners=5pt, fill=PVBlueLight!45, draw=PVBorder, dashed,
        fit=(config)(train)(valid)(best), inner sep=4mm,
        label={[font=\scriptsize\bfseries, text=PVBlue]above left:训练阶段}] {};
      \node[rounded corners=5pt, fill=PVTealLight!45, draw=PVBorder, dashed,
        fit=(protocol)(metrics)(analysis)(report), inner sep=4mm,
        label={[font=\scriptsize\bfseries, text=PVTeal]below left:评估与发布阶段}] {};
    \end{scope}
  \end{tikzpicture}
  \caption{训练、评估与结果汇总流程}
  \label{fig:experiment-pipeline}
\end{figure}

\subsection{训练目标与优化}

模型输出十类 logits，训练采用 label smoothing 系数 $\varepsilon=0.1$ 的交叉熵。设类别数为 $C=10$，平滑目标和 batch 损失为
\[
  q_{i,k}=(1-\varepsilon)\mathbf{1}[k=y_i]+\frac{\varepsilon}{C},
  \qquad
  \mathcal L=-\frac{1}{B}\sum_{i=1}^{B}\sum_{k=1}^{C}q_{i,k}\log p_{i,k}.
\]
其中 $p_{i,k}$ 为 softmax 概率。PyTorch 自动微分计算梯度，AdamW 以初始学习率 $3\times10^{-4}$ 和 weight decay 0.05 更新参数，CosineAnnealing 在训练过程中降低学习率，梯度范数裁剪上限为 1.0。训练时启用 Dropout；验证和测试使用 \texttt{eval()} 并关闭梯度记录。

\subsection{训练设置}

表\ref{tab:training-settings}列出正式实验实际使用的训练参数。六种模型沿用同一套优化设置，结构比较时不会因为学习率或 batch size 不同而混入额外变量。

\begin{table}[H]
  \centering
  \caption{正式实验的训练与选优参数}
  \label{tab:training-settings}
  \renewcommand{\arraystretch}{1.15}
  \begin{tabularx}{\textwidth}{>{\raggedright\arraybackslash}p{0.27\textwidth}>{\raggedright\arraybackslash}p{0.25\textwidth}>{\raggedright\arraybackslash}X}
    \toprule
    参数 & 设置 & 说明 \\
    \midrule
    运行设备 & CPU，8 个计算线程 & DataLoader 使用 0 个额外 worker，避免 Windows 多进程开销 \\
    随机种子 & 42、2026、3407、12345、98765 & 每个随机种子都重新完成数据划分、模型初始化、增强采样和随机协议测试 \\
    Batch size & 训练 64，评估 256 & 验证与测试不更新梯度，可以使用更大 batch \\
    训练轮数 & 最多 12 epochs & 达到早停条件时可以提前结束 \\
    优化器 & AdamW & 初始学习率 $3\times10^{-4}$，weight decay 0.05 \\
    学习率调度 & CosineAnnealing & 在最多 12 个 epoch 内逐步降低学习率 \\
    损失函数 & Cross Entropy & label smoothing 0.1 \\
    梯度裁剪 & 范数上限 1.0 & 抑制偶发的过大更新 \\
    早停 & 最少 5 epochs，patience 3 & 验证准确率改进不超过 0.0005 时计为一轮未提升 \\
    选优协议 & Center 验证集 & 保存验证准确率最高的 \texttt{best.pt} \\
    \bottomrule
  \end{tabularx}
\end{table}

在随机种子 42 的代表性运行中，ViT-Aug 于第 7 轮达到最佳验证准确率，并在第 10 轮触发早停；其余五种配置运行到第 12 轮。训练图据此显示各次运行实际完成的轮数。

所有模型均按 Center 验证准确率选择检查点，以保持选优标准一致。该设置与原始 FashionMNIST 分类任务对应，但没有直接针对 Grid 或 Shift+Rotation 表现选优。

\subsection{统一 CPU 前向基准}

除训练日志中的单轮耗时外，项目还用 \texttt{src/benchmark\_cpu.py} 对五种不同推理结构执行统一前向基准。ViT-AbsPos 与 ViT-Aug 共用同一 Tiny ViT (CLS)，数据增强只改变训练输入，因此不重复计时。所有模型均切换到 \texttt{eval()} 并关闭梯度记录，接收同一形状的随机输入；固定 8 个 PyTorch 线程和 batch size 64，先预热 5 个 batch，再重复测量 30 个 batch。表\ref{tab:cpu-benchmark}中的延迟不含数据读取、图像变换、反向传播和优化器更新，因此回答的是“同一台 CPU 上模型前向有多快”，而不是完整训练一轮需要多久。绝对数值依赖运行机器，参数量则与硬件无关。

\begin{table}[H]
  \centering
  \caption{统一 CPU 前向基准。使用 8 个 PyTorch 线程、batch size 64，预热后重复 30 次；延迟不含数据加载。}
  \label{tab:cpu-benchmark}
  \renewcommand{\arraystretch}{1.12}
  \begin{tabular}{lrrr}
    \toprule
    Model & Parameters & Mean latency (ms/batch) & Throughput (images/s) \\
    \midrule
    MLP & 805,642 & 0.37 & 173,570.3 \\
    CNN & 1,625,866 & 16.78 & 3,814.1 \\
    Tiny ViT (CLS) & 75,146 & 4.02 & 15,933.5 \\
    ViT-MeanPool & 75,018 & 3.75 & 17,067.6 \\
    HybridConv-ViT & 291,402 & 28.15 & 2,273.2 \\
    \bottomrule
  \end{tabular}
\end{table}


\subsection{检查点与日志}

每轮训练保存 \texttt{last.pt}，验证准确率刷新时保存 \texttt{best.pt}。检查点包含模型权重、优化器、学习率调度器、当前轮次、最佳验证准确率和完整配置。最终测试重新加载 \texttt{best.pt}。

日志逐轮记录训练损失、训练准确率、验证损失、验证准确率、学习率和单轮耗时。正式重新训练开始时，当前运行的旧日志会被清空，避免多次运行的数据拼接成错误曲线；smoke test 使用独立的 \texttt{\_smoke} 目录，也不会写入正式汇总表。

\subsection{评估协议与指标}

基础评估在完整官方测试集上分别构造 Center、Small Shift、Large Shift、Rotation 和 Shift+Rotation 数据。Rotation 表示居中前景在 $[-45^\circ,45^\circ]$ 连续随机旋转；Shift+Rotation 同时随机采样位置和角度。固定位置网格则对 49 个坐标各评估一次完整测试集，固定角度扫描对七个角度分别评估完整测试集。表\ref{tab:rotation-protocols}进一步区分几种旋转指标的含义。

\begin{table}[H]
  \centering
  \caption{旋转相关指标的采样协议与含义，连续随机指标和固定角度汇总来自不同测试构造}
  \label{tab:rotation-protocols}
  \begin{tabularx}{\textwidth}{>{\raggedright\arraybackslash}p{0.28\textwidth}>{\raggedright\arraybackslash}p{0.29\textwidth}X}
    \toprule
    指标 & 角度与位置协议 & 回答的问题 \\
    \midrule
    Rotation Accuracy & 居中位置；每个测试样本在 $[-45^\circ,45^\circ]$ 连续采样一次角度 & 随机旋转分布上的总体准确率 \\
    Fixed-angle Mean & 七个固定角度分别跑完整测试集，再作宏平均 & 各方向等权时的平均水平 \\
    Fixed-angle Worst & 七个固定角度准确率的最小值 & 已扫描极端方向中的最弱表现 \\
    Shift+Rotation Accuracy & 位置与角度同时连续随机采样 & 两类扰动共同出现时的总体准确率 \\
    \bottomrule
  \end{tabularx}
\end{table}

除 Accuracy 外，项目定义
\begin{align}
  \text{Robust Drop} &= \text{Center Accuracy}-\text{Large Shift Accuracy},\\
  \text{Rotation Drop} &= \text{Center Accuracy}-\text{Rotation Accuracy}.
\end{align}
二者均以百分点表示。正值越大，说明模型从中心分布转移到扰动分布时损失越明显；小幅负值并不表示扰动提高了模型能力，而通常说明增强模型在不同条件下的性能接近，随机测试采样和训练分布差异造成了数个百分点波动。

Grid Accuracy 是 49 个位置准确率的宏平均，同时报告网格最小值和最大值以分析空间平坦性。逐类准确率和混淆矩阵基于 Large Shift 预测统计，逐样本结果记录 $dx$、$dy$ 和 angle，用于错误样本与可视化分析。

\subsection{可复现运行顺序}

在 Windows PowerShell 中，正式流程只保留多随机种子训练和评估。随机种子 42 本身已包含在五个种子中，训练曲线、位置热力图、逐类结果和 Attention 对比直接从该次运行生成，无需再训练一套内容重复的单种子模型。

\begin{verbatim}
conda activate pv-vit
python src/visualize_data.py
.\scripts\run_multiseed_experiments.ps1
.\scripts\eval_multiseed_grid.ps1
python src/visualize_attention_comparison.py
.\scripts\run_augmentation_ablation.ps1
.\scripts\run_position_encoding_ablation.ps1
.\scripts\run_matched_clipping.ps1
\end{verbatim}

\texttt{eval\_multiseed\_grid.ps1} 为 30 个检查点完成固定位置和固定角度评估，再用随机种子 42 生成细节图，并汇总五个随机种子的主表和统计图。表格数据来自每次运行的 \texttt{evaluation.json}、网格 CSV 和角度 CSV。

后三个脚本分别运行三个随机种子的逐项增强消融、固定 Center+CLS 条件下的位置编码消融，以及按相同欧氏位移半径分组的边缘裁剪分析。脚本会跳过已有 \texttt{evaluation.json} 的训练结果；从头复跑前两个消融时可添加 \texttt{-Force}。
